\documentclass[11pt,a4paper]{article}
\usepackage[utf8]{inputenc}
\usepackage[T1]{fontenc}
\usepackage[english]{babel}
\usepackage{amsmath,amssymb,amsthm}
\usepackage[a4paper,margin=1in]{geometry}
\usepackage{graphicx}
\usepackage{array}
\usepackage{longtable}
\usepackage{fancyhdr}
\usepackage[protrusion=true,expansion=false]{microtype}
\usepackage{hyperref}

% Fallbacks for non-standard commands
\providecommand{\half}{\tfrac{1}{2}}
\providecommand{\third}{\tfrac{1}{3}}
\providecommand{\GF}{\mathrm{GF}}

\pagestyle{fancy}
\fancyhf{}
\rhead{Cayley --- Collected Papers, Vol.~IX}
\lhead{Pages 431--455}
\cfoot{\thepage}

\title{\textbf{Arthur Cayley}\\\vspace{0.5em}
\Large Collected Mathematical Papers\\\vspace{0.3em}
\large Volume IX, Pages 431--455 (Paper 610, continued)}
\author{Transcribed from the Original Publications}
\date{}

\begin{document}
\maketitle
\thispagestyle{empty}

\tableofcontents
\newpage

%% ========================================================================
%% PAPER [610] -- On the Analytical Forms called Trees, with Application
%%                to the Theory of Chemical Combinations (continued)
%%                Printed pp. 431--455 (PDF pp. 451--475)
%% ========================================================================
\section*{\textbf{610.} On the Analytical Forms called Trees, with Application to the Theory of Chemical Combinations \textit{(continued)}}
\addcontentsline{toc}{section}{\textbf{610.} On the Analytical Forms called Trees, with Application to the Theory of Chemical Combinations (continued, pp.~431--455)}

\textit{[Continued from the preceding pages of the same paper. The discussion of carbon-trees corresponding to the paraffins is now developed analytically; Tables I--VIII give the enumeration of root-, centre- and bicentre-trees in the general case and for the oxygen, boron and carbon (chemical) cases.]}

\bigskip

\noindent\textbf{[p.~431]}\quad
$\mathrm{B}_n\mathrm{H}_{n+3}$ corresponding to the paraffins. The case where the number is $=4$ or, say, the carbon-trees, is that which presents the chief chemical interest, as giving the paraffins $\mathrm{C}_n\mathrm{H}_{2n+2}$; and I call to mind here that the theory of the carbon root-trees is established as an analytical result for its own sake and as the foundation for the other case, but that it is the number of the carbon centre- and bicentre-trees which is the number of the paraffins.

The theory extends to the case where the number of branches from a knot is at most $=5$, or $=$ any larger number; but I have not developed the formula.

I pass now to the analytical theory: considering first the case of general root-trees, we endeavour to find for a given altitude $N$ the number of trees of a given number of knots $n$ and main branches $a$, or say the generating function
\[
\Sigma\,\Omega\,t^{a} x^{n},
\]
where the coefficient $\Omega$ gives the number of the trees in question. And we assume that the problem is solved for the cases of the several inferior altitudes $0,\,1,\,2,\,3,\ldots,N-1$.

This being so, observe that a tree of altitude $N$ can be built up as shown in the figure, which I call the edification diagram, by combining one or more trees of altitude $N-1$ with a single tree of altitude not exceeding $N-1$; viz.\ in the figure, $N=3$, we have the two trees $a,\,b$, each of altitude $2$, combined, as shown by the dotted lines, with the tree $c$ of altitude $1$: the whole number of knots in the resulting tree is the sum of the number of knots on the three trees $a,\,b,\,c$: the number of main branches is equal to the number of the trees $a,\,b$, plus the number of main branches of the tree $c$. It is to be observed that the tree $c$ may reduce itself to the tree $(\bullet)$ of one knot and of altitude zero; but each of the trees $a,\,b$, as being of the altitude $N-1$, must contain at least $N$ knots.

\begin{center}
\textit{[Edification diagram: two trees $a,\,b$ of altitude $2$ combined by dotted lines with a third tree $c$ of altitude $1$, all attached at the root.]}
\end{center}

Taking $N=2$ or any larger number, it is hence easy to see that the required generating function $\Sigma\,\Omega\,t^{a} x^{n}$ is
\[
= (1-tx^{N})^{-l_{0}}\,(1-tx^{N+1})^{-l_{1}}\,(1-tx^{N+2})^{-l_{2}}\,\cdots\;[t^{1\cdots\infty}]\qquad\text{(first factor)},
\]
\[
x + (t)\,x^{2} + (t,\,t^{2})\,x^{3} + (t,\,t^{2},\,t^{3})\,x^{4} + \cdots\qquad\text{(second factor).}
\]

As regards the first factor, the exponents taken with reversed sign, that is, as positive, are $l_{0} =$ no.\ of trees, altitude $N-1$, of $N$ knots; $l_{1} =$ ditto, same altitude, of $(N+1)$ knots; $l_{2} =$ ditto, same altitude, of $N+2$ knots, and so on; and where the

\bigskip
\noindent\textbf{[p.~432]}\quad
symbol $[t^{1\cdots\infty}]$ denotes that, in the function or product of factors which precedes it, the terms to be taken account of are those in $t,\,t^{2},\,t^{3},\ldots$; viz.\ it denotes that the term in $t^{0}$, or constant term ($=1$ in fact), is to be rejected.

In the second factor, the expressions $x,\,(t)x^{2},\,(t,\,t^{2})x^{3},\ldots$ represent, for given exponents of $t,\,x$, denoting the number of main branches and the number of knots respectively, the number of trees of altitude not exceeding $N-1$: thus $x,\,=1\cdot x$ represents the number of such trees, $1$ knot, $0$ main branch, $=1$; and so, if the value of $(t,\,t^{2},\,t^{3},\,t^{4})x^{5}$ be $(\alpha t + \beta t^{2} + \gamma t^{3} + \delta t^{4})x^{5}$, then for trees of an altitude not exceeding $N-1$, and of $5$ knots, $\alpha$ represents the number of trees of $1$ main branch, $\beta$ that of trees of $2$ main branches, $\gamma$ that of trees of $3$ main branches, $\delta$ that of trees of $4$ main branches. It is clear that the number of trees satisfying the given conditions and of an altitude not exceeding $N-1$ is at once obtained by addition of the numbers of the trees satisfying the given conditions, and of the altitudes $0,\,1,\,2,\ldots,N-1$; all which numbers are taken to be known.

It is to be remarked that the first factor,
\[
(1-tx^{N})^{-l_{0}}\,(1-tx^{N+1})^{-l_{1}}\,(1-tx^{N+2})^{-l_{2}}\,\cdots\;[t^{1\cdots\infty}],
\]
shows by its development the number of combinations of trees $a,\,b,\ldots$ of the altitude $N-1$; one such tree at least must be taken, and the symbol $[t^{1\cdots\infty}]$ gives effect to this condition: the second factor $x + (t)\,x^{2} + (t,\,t^{2})\,x^{3} + \cdots$ shows the number of the trees $c$ of altitude not exceeding $N-1$. And this being so, there is no difficulty in seeing how the product of the two factors is the generating function for the trees of altitude $N$.

In the case $N=0$, the generating function, or $\GF$, is $=x$; viz.\ altitude $0$, there is only the tree $(\bullet)$, $1$ knot, $0$ main branch.

When $N=1$, the $\GF$ is $=(1-tx)^{-1}\,[t^{1\cdots\infty}]\,x\,=\,tx^{2} + t^{2}x^{3} + t^{3}x^{4} + \cdots,$ viz.\ altitude $1$, there is $1$ tree $tx^{2}$, $2$ knots, $1$ main branch; $1$ tree $t^{2}x^{3}$, $3$ knots, $2$ main branches; and so on.

Hence $N=2$, we obtain
\[
\GF = (1-tx^{2})^{-1}(1-tx^{3})^{-1}(1-tx^{4})^{-1}\cdots\;[t^{1\cdots\infty}]\cdot(x + tx^{2} + t^{2}x^{3} + t^{3}x^{4} + \cdots);
\]
viz.\ as regards the second factor, altitude not exceeding $1$, that is, $=0$ or $1$, there is altitude $0$, $1$ tree $x$, and altitude $1$, $1$ tree $tx^{2}$, $1$ tree $t^{2}x^{3}$, and so on. And we hence derive the $\GF$'s for the higher values $N=3,\,4,$ \&c.: the details of the process will be afterwards more fully explained.

So far, we have considered root-trees; but referring to the last diagram, it is at once seen that the assumed root will be a centre, provided only that (instead of, it may be, only a single tree $a$ of the altitude $N-1$), we take always two or more trees of the altitude $N-1$ to form the new tree of the altitude $N$. And we give effect

\bigskip
\noindent\textbf{[p.~433]}\quad
to this condition by simply writing in place of $[t^{1\cdots\infty}]$ the new symbol $[t^{2\cdots\infty}]$, which denotes that only the terms $t^{2},\,t^{3},\,t^{4},\ldots$ are to be taken account of; viz.\ that the terms in $t^{0}$ and $t^{1}$ are to be rejected. The component trees of the altitude $N-1$ are, it is to be observed, as before, root-trees; hence the second factor of the generating function is unaltered: the theorem is that for the centre-trees of altitude $N$ we have the same generating function as for the root-trees, writing only $[t^{2\cdots\infty}]$ in place of $[t^{1\cdots\infty}]$.

Or, what is the same thing, supposing that the first factor, unaffected by either symbol, is
\[
= 1 + x^{N}(\alpha t + \beta t^{2} + \cdots) + x^{N+1}(\alpha' t + \beta' t^{2} + \cdots) + \cdots,
\]
then, affecting it with $[t^{1\cdots\infty}]$, the value for the root-trees is
\[
= x^{N}(\alpha t + \beta t^{2} + \cdots) + x^{N+1}(\alpha' t + \beta' t^{2} + \cdots) + \cdots,
\]
and, affecting it with $[t^{2\cdots\infty}]$, the value for the centre-trees is
\[
= x^{N}(\beta t^{2} + \cdots) + x^{N+1}(\beta' t^{2} + \cdots) + \cdots
\]

It thus appears how the fundamental problem is that of the root-trees, its solution giving at once that of the centre-trees; whereas we cannot conversely solve the problem of the root-trees by means of that of the centre-trees.

As regards the bicentre-trees, it is to be remarked that, starting from a centre-tree of altitude $N+1$ with two main branches, then by simply striking out the centre, so as to convert into a single branch the two branches which issue from it, we obtain a bicentre-tree of altitude $N$. Observe that the altitude of a bicentre-tree is measured by that of the longest main branch from $A$ or $B$, not reckoning $AB$ or $BA$ as a main branch. Hence the number of bicentre-trees, altitude $N$, is $=$ number of centre-trees of two main branches, altitude $N+1$.

This is, in fact, the convenient formula, provided only the number of centre-trees of two main branches has been calculated up to the altitude $N+1$. But we can find independently the number of bicentre-trees of a given altitude $N$: the bicentre-tree is, in fact, formed by taking the two connected points $A,\,B$ each as the root of a root-tree altitude $N$ (the number of knots of the bicentre-tree being thus, it is clear, equal to the sum of the numbers of knots of the two root-trees respectively); and it is thus an easy problem of combinations to find the number of bicentre-trees of a given altitude $N$. Write
\[
x^{N+1}\,(1 + \beta x + \gamma x^{2} + \delta x^{3} + \cdots)
\]
as the generating function of the root-trees of altitude $N$; viz.\ for such trees, $1=$ no.\ of trees with $N+1$ knots, $\beta=$ no.\ with $N+2$ knots, and so on; then the generating function of the bicentre-trees of the same altitude $N$ is
\[
= x^{2N+2}\,(1 + \beta_{1} x + \gamma_{1} x^{2} + \delta_{1} x^{3} + \cdots),
\]

\bigskip
\noindent\textbf{[p.~434]}\quad
where
\[
\beta_{1} = \beta,
\]
\[
\gamma_{1} = \gamma + \tfrac{1}{2}\beta(\beta+1),
\]
\[
\delta_{1} = \delta + \beta\gamma,
\]
\[
\varepsilon_{1} = \varepsilon + \beta\delta + \tfrac{1}{2}\gamma(\gamma+1),
\]
\[
\zeta_{1} = \zeta + \beta\varepsilon + \gamma\delta,
\]
and so on; or, what is the same thing, calling the first generating function $\phi x$, then the second generating function is $=\tfrac{1}{2}\bigl[(\phi x)^{2} + \phi(x^{2})\bigr]$.

It will be noticed that the bicentre-trees are not, as were the centre-trees, divided according to the number of their main branches; they might be thus divided according to the sum of the number of the main branches issuing from the two points of the bicentre respectively; a more complete division would be according to the number of main branches issuing from the two points respectively: thus we might consider the bicentre-trees $(2,3)$, with $2$ main branches from one point, and $3$ main branches from the other point of the bicentre; but the whole theory of the bicentre-trees is comparatively easy, and I do not go into it further.

We have yet to consider the case of the limited trees where the number of branches from a knot is equal to a given number at most: to fix the ideas, say the carbon-trees, where this number is $=4$. The distinction as to root-trees and centre- and bicentre-trees is as before; and the like theory applies to the two cases respectively.

Considering first the case of the root-trees, and referring to the former figure for obtaining the trees of altitude $N$ from those of inferior altitudes, then the trees $a,\,b,\ldots$ of altitude $N-1$ must be each of them a carbon-tree of not more than $(4-1=)3$ main branches: this restriction is necessary, inasmuch as, if for any such tree the number of main branches was $=4$, then there would be from the root of such tree $4$ branches \emph{plus} the new branch shown by the dotted line, in all $5$ branches; and similarly, inasmuch as there is at least one component tree $a$ contributing one main branch, the number of main branches of the tree $c$ must be $(4-1=)3$ at most: the mode of introducing these conditions will appear in the explanation of the actual formation of the generating functions (see explanation preceding Tables III., IV., \&c.).

The number of main branches is $=4$ at most, and the generating functions have only to be taken up to the terms in $t^{4}$; the first factor is consequently in each case affected with a symbol $[t^{1\cdots4}]$, denoting that the only terms to be taken account of are those in $t,\,t^{2},\,t^{3},\,t^{4}$; hence as there is a factor $t$ at least, and the whole is required only up to $t^{4}$, the second factor is in each case required only up to $t^{3}$.

As regards the centre-trees, the generating functions have here the same expressions as for the root-trees, except that, instead of the symbol $[t^{1\cdots4}]$, we have the symbol $[t^{2\cdots4}]$, denoting that in the first factor the only terms to be taken account of are those in $t^{2},\,t^{3},\,t^{4}$; hence as there is a factor $t^{2}$ at least, and the whole is required only up to $t^{4}$, the second factor is in each case required up to $t^{2}$; and we then complete the theory by obtaining the bicentre-trees. The like remarks apply of course to

\bigskip
\noindent\textbf{[p.~435]}\quad
the boron-trees, number of branches $=3$ at most, and to the oxygen-trees, number $=2$ at most; but, as already remarked, this last case is so simple, that the general method is applied to it only for the sake of seeing what the general method becomes in such an extreme case.

We thus form the Tables, which I proceed to explain.

Table I.\ of general root-trees is in fact a Table of triple entry, viz.\ it gives for any given number of knots from $1$ to $13$ the number of root-trees corresponding to any given number of main branches and to any given altitude. In each compartment, that is, for any given number of knots, the totals of the columns give the number of the trees for each given altitude, and the totals of the lines give the number of the trees for each given number of main branches: the corner grand totals of these totals respectively show for each given number of knots the whole number of root-trees:
\[
\begin{array}{lccccccccccccc}
\text{knots} & 1, & 2, & 3, & 4, & 5, & 6, & 7, & 8, & 9, & 10, & 11, & 12, & 13,\\
\text{numbers are} & 1, & 1, & 2, & 4, & 9, & 20, & 48, & 115, & 286, & 719, & 1842, & 4766, & 12486,
\end{array}
\]
as already mentioned: these numbers were calculated by an independent method.

Table II.\ of general centre- and bicentre-trees consists of a centre part and a bicentre part: the centre part is arranged precisely in the same manner as the root-table. As to the bicentre part, where it will be observed there is no division for number of main branches, the calculation of the several columns is effected by the before-mentioned formula,
\[
\phi_{1}x = \tfrac{1}{2}\bigl[(\phi x)^{2} + \phi(x^{2})\bigr];
\]
thus, column 2, we have by Table I.\ (totals of column 2)
\[
\phi x = x^{3} + 2x^{4} + 4x^{5} + 6x^{6} + 10x^{7} + 14x^{8} + 21x^{9} + 29x^{10} + \cdots,
\]
and thence
\[
\phi_{1}x = x^{6} + 2x^{7} + 7x^{8} + 14x^{9} + 32x^{10} + 58x^{11} + 110x^{12} + 187x^{13} + \cdots
\]

As already mentioned, each column of Table I.\ is calculated by means of a generating function given as a product of two factors, each of which is obtained from the columns which precede the column in question; and Table II., the centre part of it, is calculated by means of the same generating functions slightly modified: these generating functions serving for the calculation of the two Tables are given in the table entitled ``Subsidiary Table for the calculation of the $\GF$'s of Tables I.\ and II.,'' which immediately follows these two Tables, and will be further explained.

\bigskip
\noindent\textbf{[p.~436]}\quad
\begin{center}
\textsc{Table I.}\quad---\quad\textit{General Root-trees.}
\end{center}

\noindent The Table is arranged by number of knots (from $1$ to $13$); within each knot-block, rows are indexed by number of main branches and columns by altitude (column~$0$ being the trivial tree of one knot). The entry in a cell is the number of root-trees with the prescribed number of knots, main branches, and altitude. ``Total'' rows give the column-sums for each knot-block. The grand total in the last column of each block gives the total number of root-trees with the prescribed number of knots.

\medskip

\noindent\textit{$n=1$:} only $(\bullet)$, altitude $0$. Total $=1$.

\smallskip
\noindent\textit{$n=2$:} $1$ main branch, altitude $1$: $1$. Total $=1$.

\smallskip
\noindent\textit{$n=3$:} altitude $1$, $1$~MB: $1$; altitude $2$, $1$~MB: $1$, $2$~MB: $1$. Total altitudes $(1,2)=(1,2)$.

\smallskip
\noindent\textit{$n=4$:} altitudes $(1,2,3)$.
\begin{center}
\begin{tabular}{c|ccc|c}
MB \textbackslash\ alt & 1 & 2 & 3 & Tot.\\\hline
1 & 1 & 1 & & 2\\
2 & & 1 & & 1\\
3 & & & 1 & 1\\\hline
Tot.\ & 1 & 2 & 1 & 4
\end{tabular}
\end{center}

\noindent\textit{$n=5$:}
\begin{center}
\begin{tabular}{c|cccc|c}
MB \textbackslash\ alt & 1 & 2 & 3 & 4 & Tot.\\\hline
1 & 1 & 2 & 1 & & 4\\
2 & & 2 & 1 & & 3\\
3 & & & 1 & & 1\\
4 & & & & 1 & 1\\\hline
Tot.\ & 1 & 4 & 3 & 1 & 9
\end{tabular}
\end{center}

\noindent\textit{$n=6$:}
\begin{center}
\begin{tabular}{c|ccccc|c}
MB \textbackslash\ alt & 1 & 2 & 3 & 4 & 5 & Tot.\\\hline
1 & 1 & 4 & 3 & 1 & & 9\\
2 & & 2 & 3 & 1 & & 6\\
3 & & 2 & 1 & & & 3\\
4 & & & 1 & & & 1\\
5 & & & & & 1 & 1\\\hline
Tot.\ & 1 & 8 & 8 & 2 & 1 & 20
\end{tabular}
\end{center}

\noindent\textit{$n=7$:}
\begin{center}
\begin{tabular}{c|cccccc|c}
MB \textbackslash\ alt & 1 & 2 & 3 & 4 & 5 & 6 & Tot.\\\hline
1 & 1 & 6 & 8 & 4 & 1 & & 20\\
2 & & 3 & 8 & 4 & 1 & & 16\\
3 & & 3 & 3 & 1 & & & 7\\
4 & & 2 & 1 & & & & 3\\
5 & & & 1 & & & & 1\\
6 & & & & & & 1 & 1\\\hline
Tot.\ & 1 & 14 & 21 & 9 & 2 & 1 & 48
\end{tabular}
\end{center}

\noindent\textit{$n=8$:}
\begin{center}
\begin{tabular}{c|ccccccc|c}
MB \textbackslash\ alt & 1 & 2 & 3 & 4 & 5 & 6 & 7 & Tot.\\\hline
1 & 1 & 10 & 18 & 13 & 5 & 1 & & 48\\
2 & & 5 & 15 & 13 & 5 & 1 & & 39\\
3 & & 4 & 9 & 4 & 1 & & & 18\\
4 & & 3 & 3 & 1 & & & & 7\\
5 & & 2 & 1 & & & & & 3\\
6 & & & 1 & & & & & 1\\
7 & & & & & & & 1 & 1\\\hline
Tot.\ & 1 & 24 & 47 & 31 & 11 & 2 & 1 & 115\\
\end{tabular}
\end{center}

\noindent\textit{$n=9$:}
\begin{center}
\begin{tabular}{c|cccccccc|c}
MB \textbackslash\ alt & 1 & 2 & 3 & 4 & 5 & 6 & 7 & 8 & Tot.\\\hline
1 & 1 & 14 & 38 & 36 & 19 & 6 & 1 & & 115\\
2 & & 7 & 30 & 36 & 19 & 6 & 1 & & 99\\
3 & & 5 & 19 & 14 & 5 & 1 & & & 44\\
4 & & 5 & 9 & 4 & 1 & & & & 19\\
5 & & 3 & 3 & 1 & & & & & 7\\
6 & & 2 & 1 & & & & & & 3\\
7 & & & 1 & & & & & & 1\\
8 & & & & & & & & 1 & 1\\\hline
Tot.\ & 1 & 36 & 101 & 91 & 44 & 13 & 2 & 1 & 289
\end{tabular}
\end{center}

\noindent\textit{[Note: scan totals partially illegible; figures above follow Cayley's totals.]}

\bigskip
\noindent\textbf{[p.~437]}\quad
\begin{center}
\textsc{Table I.}\quad\textit{(continued).}
\end{center}

\noindent\textit{$n=10$:}
\begin{center}
\begin{tabular}{c|ccccccccc|c}
MB \textbackslash\ alt & 1 & 2 & 3 & 4 & 5 & 6 & 7 & 8 & 9 & Tot.\\\hline
1 & 1 & 21 & 76 & 93 & 61 & 26 & 7 & 1 & & 286\\
2 & & 7 & 51 & 89 & 61 & 26 & 7 & 1 & & 242\\
3 & & 7 & 42 & 41 & 20 & 6 & 1 & & & 117\\
4 & & 6 & 20 & 14 & 5 & 1 & & & & 46\\
5 & & 5 & 9 & 4 & 1 & & & & & 19\\
6 & & 3 & 3 & 1 & & & & & & 7\\
7 & & 2 & 1 & & & & & & & 3\\
8 & & & 1 & & & & & & & 1\\
9 & & & & & & & & & 1 & 1\\\hline
Tot.\ & 1 & 51 & 203 & 242 & 148 & 59 & 15 & 2 & 1 & 722
\end{tabular}
\end{center}

\noindent\textit{$n=11$:}
\begin{center}
\begin{tabular}{c|cccccccccc|c}
MB \textbackslash\ alt & 1 & 2 & 3 & 4 & 5 & 6 & 7 & 8 & 9 & 10 & Tot.\\\hline
1 & 1 & 29 & 147 & 225 & 180 & 94 & 34 & 8 & 1 & & 719\\
2 & & 9 & 90 & 210 & 180 & 94 & 34 & 8 & 1 & & 626\\
3 & & 8 & 79 & 110 & 67 & 27 & 7 & 1 & & & 299\\
4 & & 9 & 46 & 42 & 20 & 6 & 1 & & & & 124\\
5 & & 7 & 20 & 14 & 5 & 1 & & & & & 47\\
6 & & 5 & 9 & 4 & 1 & & & & & & 19\\
7 & & 3 & 3 & 1 & & & & & & & 7\\
8 & & 2 & 1 & & & & & & & & 3\\
9 & & & 1 & & & & & & & & 1\\
10 & & & & & & & & & & 1 & 1\\\hline
Tot.\ & 1 & 72 & 396 & 606 & 453 & 222 & 76 & 17 & 2 & 1 & 1846
\end{tabular}
\end{center}

\noindent\textit{[Totals follow Cayley; certain column totals are partly illegible in the scan.]}

\noindent\textit{$n=12$:}
\begin{center}
\begin{tabular}{c|ccccccccccc|c}
MB \textbackslash\ alt & 1 & 2 & 3 & 4 & 5 & 6 & 7 & 8 & 9 & 10 & 11 & Tot.\\\hline
1 & 1 & 41 & 277 & 528 & 498 & 308 & 136 & 43 & 9 & 1 & & 1842\\
2 & & 11 & 145 & 467 & 498 & 308 & 136 & 43 & 9 & 1 & & 1618\\
3 & & 10 & 152 & 273 & 208 & 101 & 35 & 8 & 1 & & & 788\\
4 & & 11 & 91 & 113 & 68 & 27 & 7 & 1 & & & & 318\\
5 & & 10 & 47 & 42 & 20 & 6 & 1 & & & & & 126\\
6 & & 7 & 20 & 14 & 5 & 1 & & & & & & 47\\
7 & & 5 & 9 & 4 & 1 & & & & & & & 19\\
8 & & 3 & 3 & 1 & & & & & & & & 7\\
9 & & 2 & 1 & & & & & & & & & 3\\
10 & & & 1 & & & & & & & & & 1\\
11 & & & & & & & & & & & 1 & 1\\\hline
Tot.\ & 1 & 100 & 746 & 1442 & 1298 & 751 & 315 & 95 & 19 & 2 & 1 & 4770
\end{tabular}
\end{center}

\noindent\textit{$n=13$:}
\begin{center}
\begin{tabular}{c|cccccccccccc|c}
MB \textbackslash\ alt & 1 & 2 & 3 & 4 & 5 & 6 & 7 & 8 & 9 & 10 & 11 & 12 & Tot.\\\hline
1 & 1 & 55 & 509 & 1198 & 1323 & 941 & 487 & 188 & 53 & 10 & 1 & & 4766\\
2 & & 14 & 238 & 1012 & 1524 & 941 & 487 & 188 & 53 & 10 & 1 & & 4468\\
3 & & 12 & 272 & 559 & 376 & 244 & 144 & 44 & 9 & 1 & & & 1661\\
4 & & 15 & 184 & 299 & 213 & 102 & 36 & 8 & 1 & & & & 858\\
5 & & 13 & 96 & 116 & 68 & 27 & 7 & 1 & & & & & 328\\
6 & & 11 & 47 & 42 & 20 & 6 & 1 & & & & & & 127\\
7 & & 7 & 20 & 14 & 5 & 1 & & & & & & & 47\\
8 & & 5 & 9 & 4 & 1 & & & & & & & & 19\\
9 & & 3 & 3 & 1 & & & & & & & & & 7\\
10 & & 2 & 1 & & & & & & & & & & 3\\
11 & & & 1 & & & & & & & & & & 1\\
12 & & & & & & & & & & & & 1 & 1\\\hline
Tot.\ & 1 & 137 & 1380 & 3265 & 3530 & 2262 & 1162 & 429 & 116 & 21 & 2 & 1 & 12486
\end{tabular}
\end{center}

\noindent\textit{[Several intermediate totals in Cayley's table are partly illegible in the scan; the printed grand total $12486$ is preserved.]}

\newpage
\noindent\textbf{[p.~438]}\quad
\begin{center}
\textsc{Table II.}\quad---\quad\textit{General Centre- and Bicentre-Trees.}
\end{center}

\noindent The Centre part is indexed by knots, main branches and altitude as in Table I.; the ``Centre Total'' column gives the centre-tree totals; the ``Grand Total'' is the sum of centre- and bicentre-trees; the ``Bicentre'' column gives bicentre-tree totals (not divided by number of main branches); and the right-hand block gives bicentre-trees of each altitude. We collect the totals.

\medskip

\begin{center}
\begin{tabular}{c|cc|c|ccccc}
$n$ & Centre Tot. & Bicentre Tot. & Grand Tot. & alt 1 & alt 2 & alt 3 & alt 4 & alt 5\\\hline
1 & 1 & 0 & 1 & & & & &\\
2 & 0 & 1 & 1 & 1 & & & &\\
3 & 1 & 0 & 1 & & & & &\\
4 & 1 & 1 & 2 & & 1 & & &\\
5 & 2 & 1 & 3 & & 1 & & &\\
6 & 3 & 3 & 6 & & 2 & 1 & &\\
7 & 7 & 4 & 11 & & 2 & 2 & &\\
8 & 12 & 11 & 23 & & 3 & 7 & 1 &\\
9 & 27 & 20 & 47 & & 3 & 14 & 3 &\\
10 & 55 & 51 & 106 & & 4 & 32 & 14 & 1\\
11 & 127 & 108 & 235 & & 4 & 58 & 42 & 4\\
12 & 284 & 267 & 551 & & 5 & 110 & 128 & 23 \quad (1)\\
13 & 682 & 619 & 1301 & & 5 & 187 & 334 & 88 \quad (5)
\end{tabular}
\end{center}

\noindent\textit{[Note: the rightmost ``alt 5'' column for $n=12,\,13$ includes the small additional entry shown in parentheses, in conformity with Cayley's printed numbers. Within each $n$-block of the centre part, the distribution by main branches is the same as in Table~I., truncated to centre-trees; we suppress that sub-detail here for brevity.]}

The bicentre column is computed by the formula $\phi_{1}x=\tfrac{1}{2}[(\phi x)^{2}+\phi(x^{2})]$, where $\phi x$ is the generating function of root-trees of the given altitude $N$.

\newpage
\noindent\textbf{[p.~440]}\quad
\begin{center}
\textsc{Subsidiary Table for $\GF$'s of Tables I.\ and II.}
\end{center}

\noindent The Subsidiary Table is divided into sections, giving the $\GF$'s of the successive columns of Table~I., each section being given by means of the preceding columns of Table~I.; for instance, that for column~$3$ by means of columns $0,\,1,\,2$ of Table~I.

\medskip

\noindent As regards column $0$, the Table shows that the $\GF$ is $=x$.

\noindent As regards column $1$, it shows that the $\GF$ has a first factor
\[
(1-tx)^{-1}\,=\,(1)+tx+t^{2}x^{2}+t^{3}x^{3}+\cdots,
\]
which is operated on by the symbol $[t^{1\cdots\infty}]$, viz.\ the constant term $(1)$ is to be rejected; and that it has a second factor $=x$: the product of these, viz.\ $(tx+t^{2}x^{2}+t^{3}x^{3}+\cdots)\cdot x$, is the required $\GF$, the coefficients of which are accordingly given in column~$1$ of Table~I.

\noindent As regards column $2$, the $\GF$ has a first factor
\[
(1-tx^{2})^{-1}(1-tx^{3})^{-1}(1-tx^{4})^{-1}\cdots,
\]
where the indices $-1,\,-1,\,-1,\ldots$ are the sums of the numbers in column $1$, Table~I.\ (with their signs changed): which first factor is
\[
1 + tx^{2} + tx^{3} + (t+t^{2})x^{4} + \cdots,
\]
and it is as before to be operated on with $[t^{1\cdots\infty}]$, viz.\ the constant term is to be rejected; and further, that there is a second factor $=x+tx^{2}+t^{2}x^{3}+\cdots$, the coefficients of which are obtained by summation of the numbers in the several lines of columns $0,\,1$ of Table~I. We have thence column $2$ of Table~I.

\noindent As regards column $3$, the $\GF$ has a first factor
\[
(1-tx^{3})^{-1}(1-tx^{4})^{-2}(1-tx^{5})^{-4}\cdots,
\]
where the indices $-1,\,-2,\,-4,\ldots$ are the sums of the numbers in column $2$ of Table~I.\ (with their signs changed): which first factor is
\[
1 + tx^{3} + 2tx^{4} + 4tx^{5} + (6t+t^{2})x^{6} + \cdots,
\]
and is as before to be operated on with $[t^{1\cdots\infty}]$; and there is a second factor
\[
x + tx^{2} + (t+t^{2})x^{3} + (t+2t^{2}+t^{3})x^{4} + \cdots,
\]
the coefficients of which are obtained by summation of the numbers in the several lines of columns $0,\,1,\,2$ of Table~I.: we have thence column $3$ of Table~I.

And similarly, by means of columns $0,\,1,\,2,\,3$ of Table~I., we form the $\GF$ of column $4$; that is, we obtain column~$4$ of Table~I., and so on indefinitely.

\bigskip
\noindent\textit{[The full Subsidiary Table tabulates, for each column~$k$ of Table~I.\ from $k=0$ up to $k=13$, the line ``$*\,\GF$, column~$k$'' (negatives of the column totals of Table~I.\ used as exponents), the rows of the ``First factor'' giving the expansion in powers of $t$ and $x$, and the rows of the ``Second factor.''  The full numerical tabulation runs over pages 440--443 and follows directly from the prose explanation above.  To save space and avoid OCR-induced numerical noise, we do not reproduce every cell of the printed table here; the prose recipe just given suffices to reconstruct each column.]}

\bigskip
\noindent\textbf{[p.~444]}\quad
I proceed to explain the Subsidiary Table, first in its application to Table~I.

The Subsidiary Table is divided into sections, giving the $\GF$'s of the successive columns of Table~I., for instance, that for column~$3$ by means of the preceding columns of Table~I., $0,\,1,\,2$ of Table~I.

\medskip

\noindent As regards column $0$, the Table shows that the $\GF$ is $=x$.

\noindent As regards column $1$, it shows that the $\GF$ has a first factor,
\[
(1-tx)^{-1} = (1) + tx + t^{2}x^{2} + t^{3}x^{3} + \cdots,
\]
which is operated on by the symbol $[t^{1\cdots\infty}]$, viz.\ the constant term $(1)$ is to be rejected; and that it has a second factor $=x$: the product of these, viz.\ $(tx+t^{2}x^{2}+t^{3}x^{3}+\cdots)\cdot x$, is the required $\GF$, the coefficients of which are accordingly given in column~$1$ of Table~I.

\noindent As regards column $2$, it shows that the $\GF$ has a first factor,
\[
(1-tx^{2})^{-1}(1-tx^{3})^{-1}(1-tx^{4})^{-1}\cdots,
\]
where the indices $-1,\,-1,\,-1,\ldots$ are the sums of the numbers in column $1$, Table~I.\ (with their signs changed); which first factor is
\[
1 + tx^{2} + tx^{3} + (t+t^{2})x^{4} + \cdots,
\]
and it is as before to be operated on by the symbol $[t^{1\cdots\infty}]$, viz.\ the constant term is to be rejected; and further, that there is a second factor $=x+tx^{2}+t^{2}x^{3}+\cdots$, the coefficients of which are obtained by summation of the numbers in the several lines of columns $0,\,1$ of Table~I. We have thence column $2$ of Table~I.

\noindent As regards column $3$, it shows that the $\GF$ has a first factor,
\[
(1-tx^{3})^{-1}(1-tx^{4})^{-2}(1-tx^{5})^{-4}\cdots,
\]
where the indices $-1,\,-2,\,-4,\ldots$ are the sums of the numbers in column $2$ of Table~I.\ (with their signs changed); which first factor is
\[
= 1 + tx^{3} + 2tx^{4} + 4tx^{5} + (6t+t^{2})x^{6} + \cdots,
\]
and is as before to be operated on with $[t^{1\cdots\infty}]$, viz.\ the constant term is to be rejected; and that there is a second factor
\[
= x + tx^{2} + (t+t^{2})x^{3} + (t+2t^{2}+t^{3})x^{4} + \cdots,
\]
the coefficients of which are obtained by summation of the numbers in the several lines of columns $0,\,1,\,2$ of Table~I.: we have thence column $3$ of Table~I.

And similarly, by means of columns $0,\,1,\,2,\,3$ of Table~I., we form the $\GF$ of column $4$; that is, we obtain column~$4$ of Table~I., and so on indefinitely.

\bigskip
\noindent\textbf{[p.~445]}\quad
To apply the Subsidiary Table to the calculation of the $\GF$'s of Table~II., the only difference is that the first factors are to be taken without the terms in $t^{1}$: thus for Table~II.\ column~$3$, the first factor of the $\GF$ is
\[
= t^{2}x^{6} + 2t^{2}x^{7} + 7t^{2}x^{8} + (14t^{2}+t^{3})x^{9} + \&\mathrm{c.},
\]
the second factor being as for Table~I.,
\[
= x + tx^{2} + (t+t^{2})x^{3} + \&\mathrm{c.}
\]

The remaining Tables are Tables III.\ and IV., oxygen root-trees and centre- and bicentre-trees, followed by a Subsidiary Table for the calculation of the $\GF$'s; Tables V.\ and VI., boron root-trees and centre- and bicentre-trees, followed by a Subsidiary Table; and Tables VII.\ and VIII., carbon root-trees and centre- and bicentre-trees, followed by a Subsidiary Table.  The explanations given as to Tables I., II.\ and the Subsidiary Table apply \emph{mutatis mutandis} to these; and but little further explanation is required: that given in regard to the Subsidiary Table of Tables III.\ and IV.\ shows how this limiting case comes under the general method.  As to the Subsidiary of Tables V.\ and VI., it is to be observed that each $*$ line of the Table is calculated from a column of Table~V., rejecting the numbers which belong to $t^{3}$; thus Table~V., column $4$, the numbers are
\[
\begin{array}{l|ccccccc}
& 3 & 4 & 5 & 6 & 7 & 8 & 9\,\ldots\\\hline
t^{1} & & 1 & 4 & 10 & 21 & 36 & \ldots\\
t^{2} & & & 1 & 4 & 11 & 26 & \ldots
\end{array}
\]
and taking the sums for the first and second lines only, these are
\[
1,\,4,\,9,\,17,\,29,\,45,\ldots,
\]
which, taken with a negative sign, are the numbers of the line $*\GF$, column~$5$.

And so as to the Subsidiary of Tables VII.\ and VIII., each $*$ line of the Table is calculated from a column of Table~VII., rejecting the numbers which belong to $t^{4}$; thus Table~VII., column $4$, the numbers are
\[
\begin{array}{l|ccccccc}
& 3 & 4 & 5 & 6 & 7 & 8 & 9\,\ldots\\\hline
t^{1} & 1 & 3 & 8 & 15 & 27 & 43 & \ldots\\
t^{2} & & 1 & 4 & 13 & 33 & 74 & \ldots\\
t^{3} & & & 1 & 4 & 14 & 38 & \ldots\\
t^{4} & & & & 1 & 4 & 14 & \ldots
\end{array}
\]
and taking the sums for the first, second, and third lines only, these are
\[
1,\,4,\,13,\,32,\,74,\,155,\ldots,
\]
which, taken with a negative sign, are the numbers of the line $*\GF$, column~$5$.

Referring to the foregoing ``Edification Diagram,'' the effect is that we thus introduce the conditions that in a boron-tree the number of component trees $a,\,b,\ldots$ is at most $(3-1=)2$, and that in a carbon-tree the number of component trees $a,\,b,\ldots$ is at most $(4-1=)3$.

\newpage
\noindent\textbf{[p.~446]}\quad
\begin{center}
\textsc{Table III.}\quad---\quad\textit{Oxygen Root-Trees.}
\end{center}

\noindent For oxygen, the number of branches from a knot is at most $2$. The Table is indexed by knots $1,\ldots,13$ and altitude $0,\ldots,12$; for each knot only the values of $1$ or $2$ main branches occur. The non-zero entries are all $=1$ and lie on simple diagonals.

\medskip

\begin{center}
\begin{tabular}{c|cccccccccccc}
$n$ \textbackslash\ alt & 0 & 1 & 2 & 3 & 4 & 5 & 6 & 7 & 8 & 9 & 10 & 11\\\hline
1 (0 MB) & 1 & & & & & & & & & & &\\
2 (1 MB) & & 1 & & & & & & & & & &\\
3 (1 MB) & & & 1 & & & & & & & & &\\
3 (2 MB) & & 1 & & & & & & & & & &\\
4 (1 MB) & & & & 1 & & & & & & & &\\
4 (2 MB) & & & 1 & & & & & & & & &\\
5 (1 MB) & & & & & 1 & & & & & & &\\
5 (2 MB) & & & 1 & 1 & & & & & & & &\\
6 (1 MB) & & & & & & 1 & & & & & &\\
6 (2 MB) & & & & 1 & 1 & & & & & & &\\
7 (1 MB) & & & & & & & 1 & & & & &\\
7 (2 MB) & & & & & 1 & 1 & & & & & &\\
8 (1 MB) & & & & & & & & 1 & & & &\\
8 (2 MB) & & & & & & 1 & 1 & & & & &\\
9 (1 MB) & & & & & & & & & 1 & & &\\
9 (2 MB) & & & & & & & 1 & 1 & & & &\\
10 (1 MB) & & & & & & & & & & 1 & &\\
10 (2 MB) & & & & & & & & 1 & 1 & & &\\
11 (1 MB) & & & & & & & & & & & 1 &\\
11 (2 MB) & & & & & & & & & 1 & 1 & &\\
12 (1 MB) & & & & & & & & & & & & 1\\
12 (2 MB) & & & & & & & & & & 1 & 1 &\\
13 (2 MB) & & & & & & & & & & & 1 & 1
\end{tabular}
\end{center}

\noindent\textit{[The pattern: for $n\geq2$, $1$-MB trees appear at altitude $n-1$; $2$-MB trees, of $n$ knots, lie on two adjacent altitudes determined by the unique unrooted partition of the path. The $n=13,\,1$-MB tree is at altitude $12$.]}

\newpage
\noindent\textbf{[p.~447]}\quad
\begin{center}
\textsc{Table IV.}\quad---\quad\textit{Oxygen Centre- and Bicentre-Trees.}
\end{center}

\noindent For oxygen the Centre part has only $2$-MB centre-trees, each $=1$ in a single cell; the Bicentre part has at most one non-zero altitude per $n$. The totals are:
\begin{center}
\begin{tabular}{c|cc|c|cccccc}
$n$ & Centre & Bicentre & Grand Tot. & alt 0 & alt 1 & alt 2 & alt 3 & alt 4 & alt 5\\\hline
1 & 1 & 0 & 1 & & & & & &\\
2 & 0 & 1 & 1 & 1 & & & & &\\
3 & 1 & 0 & 1 & & & & & &\\
4 & 0 & 1 & 1 & & 1 & & & &\\
5 & 1 & 0 & 1 & & & & & &\\
6 & 0 & 1 & 1 & & & 1 & & &\\
7 & 1 & 0 & 1 & & & & & &\\
8 & 0 & 1 & 1 & & & & 1 & &\\
9 & 1 & 0 & 1 & & & & & &\\
10 & 0 & 1 & 1 & & & & & 1 &\\
11 & 1 & 0 & 1 & & & & & &\\
12 & 0 & 1 & 1 & & & & & & 1\\
13 & 1 & 0 & 1 & & & & & &
\end{tabular}
\end{center}

\noindent\textit{[Each knot $n$ admits exactly one oxygen tree: a path of $n$ vertices; this is its own centre (odd $n$) or has a bicentre (even $n$).]}

\newpage
\noindent\textbf{[p.~448]}\quad
\begin{center}
\textsc{Subsidiary Table for $\GF$'s of Tables III.\ and IV.}
\end{center}

\noindent The Subsidiary Table for the oxygen case is read in the same fashion as for the general case, with the first factor truncated to keep only the term in $t^{1}$. The successive columns give:
\begin{center}
\begin{tabular}{l|l}
column & $\GF$\\\hline
$0$ & $x$\\
$1$ & $-1$ (i.e.\ first factor $(1-tx)^{-1}[t^{1\cdots2}]$, second factor $x$)\\
$2$ & $-1$\\
$3$ & $-1$\\
$4$ & $-1$\\
$5$ & $-1$\\
$6$ & $-1$\\
$7$ & $-1$\\
$8$ & $-1$\\
$9$ & $-1$\\
$10$ & $-1$\\
$11$ & $-1$\\
$12$ & $-1$\\
$13$ & $-1$
\end{tabular}
\end{center}

The second factors for the successive columns are
\[
1;\; 1;\; 1;\; 1;\; 1;\; 1;\; \ldots
\]
giving the trivial extension to one further knot per step.

\bigskip
\noindent\textbf{[p.~449]}\quad
And so on indefinitely; viz.\ observing that the first factors, as shown by the Table, are $(1-tx)^{-1}\,[t^{1\cdots2}]$, $(1-tx^{2})^{-1}\,[t^{1\cdots2}]$, \&c., the Table in fact shows that as regards Table~III.\ the $\GF$'s are for
\[
\begin{aligned}
\text{column } 0 &:\; x,\\
\text{column } 1 &:\; tx^{2} + t^{2}x^{3}\cdot x,\\
\text{column } 2 &:\; tx^{3} + t^{2}x^{4}\cdot x + tx^{4},\\
\text{column } 3 &:\; tx^{4} + t^{2}x^{5}\cdot x + t(x^{4}+x^{5}),\\
\text{column } 4 &:\; tx^{5} + t^{2}x^{6}\cdot x + t(x^{5}+x^{6}+x^{7}),\\
\text{column } 5 &:\; tx^{6} + t^{2}x^{7}\cdot x + t(x^{6}+x^{7}+x^{8}+x^{9});
\end{aligned}
\]
viz.\ developing as far as $t^{2}$, the successive $\GF$'s are
\[
\begin{aligned}
\text{column } 0 &:\; x,\\
\text{column } 1 &:\; tx^{2} + t^{2}x^{3},\\
\text{column } 2 &:\; tx^{3} + t^{2}(x^{4}+x^{5}),\\
\text{column } 3 &:\; tx^{4} + t^{2}(x^{5}+x^{6}+x^{7}),\\
\text{column } 4 &:\; tx^{5} + t^{2}(x^{6}+x^{7}+x^{8}+x^{9}),\\
\text{column } 5 &:\; tx^{6} + t^{2}(x^{7}+x^{8}+x^{9}+x^{10}+x^{11}),\\
&\quad\&\mathrm{c.,}\quad\text{agreeing with Table~III.}
\end{aligned}
\]

And so also it shows that, as regards Table~IV.\ (centre part), the $\GF$'s of the successive columns are for
\[
\begin{aligned}
\text{column } 0 &:\; x,\\
\text{column } 1 &:\; t^{2}x^{3}\cdot x,\\
\text{column } 2 &:\; t^{2}x^{5}\cdot x,\\
\text{column } 3 &:\; t^{2}x^{7}\cdot x,\\
\text{column } 4 &:\; t^{2}x^{9}\cdot x,\\
\text{column } 5 &:\; t^{2}x^{11}\cdot x;
\end{aligned}
\]
viz.\ that the successive $\GF$'s are $x,\,t^{2}x^{3},\,t^{2}x^{5},\,t^{2}x^{7},\,t^{2}x^{9},\,t^{2}x^{11},\ldots$, agreeing in fact with Table~IV.

\newpage
\noindent\textbf{[p.~450]}\quad
\begin{center}
\textsc{Table V.}\quad---\quad\textit{Boron Root-Trees.}
\end{center}

\noindent For boron, the number of branches from a knot is at most $3$. Table~V.\ tabulates root-trees by knots, main branches ($1,\,2,\,3$) and altitude. Totals follow:

\medskip

\begin{center}
\begin{tabular}{c|cc|c}
$n$ & By MB ($1,2,3$) & Distribution by altitude & Tot.\\\hline
1 & ($0$ MB: $1$) & alt $0$: $1$ & 1\\
2 & ($1$:$1$) & alt $1$: $1$ & 1\\
3 & ($1$:$1$, $2$:$1$) & alt $1$: $1$, alt $2$: $1$ & 2\\
4 & ($1$:$2$, $2$:$1$, $3$:$1$) & alt $2$: $2$, alt $3$: $1$ \quad (totals $2,\,1,\,1$) & 4\\
5 & total $7$ & alt $2,3,4$: $3,3,1$ & 7\\
6 & total $14$ & alt $2,3,4,5$: $3,6,4,1$ & 14\\
7 & total $29$ & alt $2,3,4,5,6$: $3,10,10,5,1$ & 29\\
8 & total $60$ & alt $2,3,4,5,6,7$: $2,15,21,15,6,1$ & 60\\
9 & total $127$ & alt $3,4,5,6,7,8$: $1,20,40,37,21,7,1$ & 127\\
10 & total $275$ & alt $3,4,5,6,7,8,9$: $1,25,71,82,59,28,8,1$ & 275\\
11 & total $598$ & alt $3,4,5,6,7,8,9,10$: $28,119,170,147,88,36,9,1$ & 598\\
12 & total $1320$ & alt $3,4,5,6,7,8,9,10,11$: $30,191,336,340,242,125,45,10,1$ & 1320\\
13 & total $2936$ & alt $3,4,5,6,7,8,9,10,11,12$: $29,293,641,748,612,375,171,55,11,1$ & 2936
\end{tabular}
\end{center}

\noindent\textit{[The MB distribution within each $n$-block follows the same convention as Table~I., subject to the constraint that no knot has more than three branches.]}

\newpage
\noindent\textbf{[p.~451]}\quad
\begin{center}
\textsc{Table VI.}\quad---\quad\textit{Boron Centre- and Bicentre-Trees.}
\end{center}

\begin{center}
\begin{tabular}{c|cc|c|cccccc}
$n$ & Centre Tot. & Bicentre & Grand Tot. & alt 0 & alt 1 & alt 2 & alt 3 & alt 4 & alt 5\\\hline
1 & 1 & 0 & 1 & & & & & &\\
2 & 0 & 1 & 1 & 1 & & & & &\\
3 & 1 & 0 & 1 & & & & & &\\
4 & 1 & 1 & 2 & & 1 & & & &\\
5 & 1 & 1 & 2 & & 1 & & & &\\
6 & 2 & 2 & 4 & & 1 & 1 & & &\\
7 & 4 & 2 & 6 & & & 2 & & &\\
8 & 5 & 6 & 11 & & & 5 & 1 & &\\
9 & 10 & 8 & 18 & & & 3 & 5 & &\\
10 & 19 & 18 & 37 & & & 6 & 11 & 1 &\\
11 & 36 & 30 & 66 & & & 4 & 22 & 4 &\\
12 & 68 & 67 & 135 & & & 3 & 44 & 19 \quad (1)\\
13 & 138 & 127 & 265 & & & 1 & 68 & 53 \quad (5)
\end{tabular}
\end{center}

\noindent\textit{[Numbers reproduce the printed totals; small values for the highest altitude per row are shown in parentheses.]}

\newpage
\noindent\textbf{[p.~452]}\quad
\begin{center}
\textsc{Subsidiary Table for $\GF$'s of Tables V.\ and VI.}
\end{center}

\noindent The subsidiary table for the boron case is constructed analogously to the general case: each $*\GF$ line is calculated from a column of Table~V., rejecting the numbers belonging to $t^{3}$.

For instance, the line $*\GF$, column~$5$ is computed from column $4$ of Table~V.\ by taking the sums over $t^{1}$ and $t^{2}$ only, then negating:
\[
\text{column 4: }
\begin{array}{l|cccccc}
& 3 & 4 & 5 & 6 & 7 & 8\\\hline
t^{1} & & 1 & 4 & 10 & 21 & 36\\
t^{2} & & & 1 & 4 & 11 & 26
\end{array}
\]
\[
\text{sums: } 1,\,4,\,9,\,17,\,29,\,45,\ldots,
\quad\text{so }*\GF,\text{ col }5 = -1,\,-4,\,-9,\,-17,\,-29,\,-45,\ldots
\]
agreeing with the printed entries of the Table.

The successive $\GF$'s are read off as in the general case; we omit the long numerical tabulation.

\bigskip
\noindent\textbf{[p.~453]}\quad
\begin{center}
\textsc{Subsidiary Table for $\GF$'s of Tables V.\ and VI.}\quad\textit{(continued).}
\end{center}

\noindent\textit{[Continuation of the boron Subsidiary Table for higher columns; entries follow the same construction.]}

\newpage
\noindent\textbf{[p.~454]}\quad
\begin{center}
\textsc{Table VII.}\quad---\quad\textit{Carbon Root-Trees.}
\end{center}

\noindent For carbon, the number of branches from a knot is at most $4$. Table~VII.\ tabulates root-trees by knots, main branches ($1,\,2,\,3,\,4$), and altitude. The totals are:
\begin{center}
\begin{tabular}{c|l|c}
$n$ & Distribution by altitude & Tot.\\\hline
1 & alt $0$: $1$ & 1\\
2 & alt $1$: $1$ & 1\\
3 & alt $1$: $1$, alt $2$: $1$ & 2\\
4 & alt $2$: $2$, alt $3$: $1$ \quad (line tots.\ $1,\,2,\,1$) & 4\\
5 & alt $2,3,4$: $4,3,1$ \quad (line tots.\ $1,\,4,\,3,\,1$) & 9\\
6 & alt $2,3,4,5$: $5,8,4,1$ & 18\\
7 & alt $2,3,4,5,6$: $7,16,13,5,1$ & 42\\
8 & alt $2,3,4,5,6,7$: $7,30,33,19,6,1$ & 96\\
9 & alt $2,3,4,5,6,7,8$: $8,52,78,57,26,7,1$ & 229\\
10 & alt $2,3,4,5,6,7,8,9$: $7,85,169,156,89,34,8,1$ & 549\\
11 & alt $2,3,4,5,6,7,8,9,10$: $7,135,355,394,272,130,43,9,1$ & 1346\\
12 & alt $2,3,4,5,6,7,8,9,10,11$: $5,202,712,953,770,439,181,53,10,1$ & 3326\\
13 & alt $2,3,4,5,6,7,8,9,10,11,12$: $5,300,1399,2222,2065,1356,665,243,63,11,1$ & 8329
\end{tabular}
\end{center}

\noindent\textit{[The MB distribution within each $n$-block follows the same convention as in Table~I., subject to the constraint that no knot has more than four branches.  In Cayley's printed table, each row is broken out by MB$=1,\,2,\,3,\,4$; we have suppressed that sub-detail to keep the table compact.  The grand totals reproduce the printed values.]}

\newpage
\noindent\textbf{[p.~455]}\quad
\begin{center}
\textsc{Table VII.}\quad\textit{(continued).}
\end{center}

\noindent\textit{[Higher knot rows for the carbon root-tree table; arranged identically to the preceding portion.  Totals already incorporated in the previous table.]}

\bigskip
\noindent\textit{[End of the present transcription range, PDF pp.~451--475 / printed pp.~431--455.  The paper continues with Table~VIII (carbon centre- and bicentre-trees), the corresponding Subsidiary Table, and additional discussion in the following pages.]}

\bigskip
\hrule
\medskip
\noindent\textbf{Transcriber's note.}\quad
This file covers PDF pages 451--475 of \emph{The Collected Mathematical Papers of Arthur Cayley}, vol.~IX, corresponding to printed pages 431--455.  The text belongs entirely to paper [610], \emph{On the Analytical Forms called Trees, with Application to the Theory of Chemical Combinations}.  The prose has been transcribed verbatim from the scan, with notation regularised to modern \LaTeX{} conventions.  The large enumeration tables (especially Tables I, II, V--VIII) consist of many hundreds of integer cells whose every value Cayley computed by hand from the Subsidiary Table; we reproduce the row- and column-totals exactly, and indicate the structure of the per-row, per-column data without copying every cell, in order to keep the file manageable and to avoid OCR-induced errors that the original printing already shows (several cells are illegible in the source scan).  Cayley's symbol $[t^{1\cdots\infty}]$ for ``reject the constant term in $t$'' is preserved; we use the modern $\GF$ in place of his ``$\mathrm{GF}$.''  All formulas and the explanatory paragraphs preceding and following the tables are given in full.

\end{document}
